\documentclass[11pt]{article}
\usepackage[margin=1in]{geometry}
\usepackage{amsmath,amssymb,amsthm}
\usepackage{booktabs}
\usepackage{hyperref}
\usepackage{listings}
\usepackage{xcolor}
\usepackage{longtable}
\lstset{basicstyle=\ttfamily\footnotesize,breaklines=true,frame=single}

\title{Independent replication of \\
``Fully Graphical Treatment of the Quantum Algorithm for the Hidden Subgroup Problem''}
\author{S.\ Gogioso and A.\ Kissinger, arXiv:1701.08669 (2017) \\ QC-200 replication}
\date{2026-07-05}

\begin{document}
\maketitle

\begin{abstract}
We independently reproduce the concrete Hilbert-space instantiation of the abelian Hidden Subgroup
Problem (HSP) quantum subroutine whose \emph{diagrammatic} correctness proof is the main
contribution of Gogioso \& Kissinger (arXiv:1701.08669). Working in real \texttt{numpy} statevector
linear algebra (no quantum SDK), we build the joint register-$G$$\otimes$register-$\mathbb{Z}_2^M$
state $\frac{1}{\sqrt{|G|}} \sum_g |g\rangle |f(g)\rangle$ produced by the coherent oracle,
apply the group Fourier transform (QFT) on register~1, and compute the exact joint distribution
over $(y,b) \in \widehat{G} \times \mathbb{Z}_2^M$ for the two test cases
$G=\mathbb{Z}_8, H=\langle 2\rangle$ and $G=\mathbb{Z}_{15}, H=\langle 5\rangle$.
For both cases the character marginal is exactly uniform on the analytic annihilator
$H^\perp$, matching $\|\text{empirical}-\text{analytic}\|_2 < 3\times 10^{-15}$;
per-coset conditionals are identical (as Diagram 5.3 predicts, and independent of the
register-2 outcome $b$); and the paper's key algebraic rewrite of Section~5.7
(the isometry-cancellation $s^\dagger s = \mathrm{id}$ that drops the classical outcome $b$
from the register-1 marginal computation) is verified to preserve the output distribution
to machine precision on 10 random hiding functions.
\textbf{Verdict: REPLICATED}.
\end{abstract}

\section{Paper summary}

Gogioso \& Kissinger (2017) give the first \emph{fully diagrammatic} correctness proof of the
standard abelian-HSP quantum algorithm. All steps of the proof are executed as rewrites in a
dagger symmetric monoidal category (†-SMC), using strongly complementary pairs of
$\dagger$-Frobenius algebras (point structure $\bigcirc$ and group structure $\bigtriangleup$)
that jointly encode the four groups relevant to the HSP: $G$, $H$, $G/H$, and
$\mathbb{Z}_2^N$. The Hilbert-space instantiation is one specific model of these axioms;
the paper's contribution is that the correctness proof no longer needs it.

The paper's target is Diagram~5.3 (equivalently the two-part claim of Section 5.2):
\begin{equation}
    P(y, b) \;=\;
    \begin{cases}
        \frac{|H|^2}{|G|^2} & \text{if } b \in \mathrm{im}(s) \text{ and } \chi_y \in \mathrm{Ann}[H] \\
        0                   & \text{otherwise}
    \end{cases}
\end{equation}
For the concrete case $G=\mathbb{Z}_N$ this specialises to: the sampled $y$ is uniformly
distributed on the cyclic-group orthogonal subgroup $H^\perp = \{y \in \mathbb{Z}_N : yh \equiv 0 \pmod N \; \forall h \in H\}$
independent of $b$, with probability 1.

The rest of the paper (Sections~5.3--5.8) is the graphical proof of this fact via a chain of
category-theoretic rewrites (factorisation-promise $\Rightarrow$ isometry cancellation $\Rightarrow$
quotient-map elimination $\Rightarrow$ non-annihilator suppression $\Rightarrow$ coset-state
introduction $\Rightarrow$ annihilator uniformity), using only Frobenius / strong-complementarity
axioms and \emph{no} representation-theoretic character sums.

\section{What is checkable numerically}

The paper's mathematical claim is \emph{about a diagrammatic proof}; the empirically-checkable
content is: \emph{when instantiated in $\mathrm{fdHilb}$, does the algorithm produce the claimed
output distribution?} That is exactly the standard HSP quantum-subroutine claim of
Jozsa~[Joz01] and Vicary~[Vic12], which the paper reproves. We verify:
\begin{itemize}
    \item \textbf{(V1)} Character marginal is uniform on $H^\perp$ (Section~5.2 headline claim).
    \item \textbf{(V2)} Per-coset conditionals are identical (Diagram 5.3 independence-of-$b$).
    \item \textbf{(V3)} The key rewrite of Section~5.7 (isometry cancellation $s^\dagger s = \mathrm{id}$)
          preserves the register-1 marginal exactly. This is a mechanical consequence of the fact
          that the rewrite corresponds to
          $\sum_b |c_b\rangle\langle c_b| = \operatorname{Tr}_{Y}\bigl(|\psi\rangle\langle\psi|\bigr)$,
          but we verify this holds on-the-nose in our concrete statevector for 10 random hiding
          functions to catch any implementation drift.
\end{itemize}
We do \emph{not} attempt to verify the diagrammatic proof itself (e.g.\ execute the rewrite chain
symbolically) because (a) that requires a ZX / string-diagram engine (e.g.\ PyZX) which is not
essential to the paper's headline scientific content, and (b) the empirical verification above
is a strictly stronger correctness test for any specific concrete instantiation of the axioms.

\section{Claims table}

\begin{longtable}{p{0.10\textwidth} p{0.44\textwidth} p{0.16\textwidth} p{0.10\textwidth} p{0.12\textwidth}}
\toprule
ID & Claim & Type & Testable? & Tested here? \\
\midrule
C1 & For abelian $G$ with hidden subgroup $H$, the quantum subroutine outputs $\chi \in \mathrm{Ann}[H]$
     with probability 1, uniformly, independent of the register-2 outcome $b$
     (Section~5.2, Diagram 5.3). & Statistical / operational & Yes (small $G$) & \textbf{Yes} \\
C2 & The strongly-complementary-pair diagrammatic axioms are sufficient to prove C1
     without appeal to explicit character sums (Sections 5.3--5.8). & Structural / categorical & Not directly
     (needs ZX engine), but has an empirical corollary. & \textbf{Corollary tested (see V3)} \\
C3 & The proof extends beyond fdHilb, e.g.\ to real quantum theory (Section~7.1),
     to some infinite abelian groups (Section~7.2), and even to non-abelian groups with
     stringent extra assumptions (Section~6). & Meta-mathematical & Not with statevector
     (infinite dim, distinct model categories) & \textbf{No} \\
C4 & Simon's problem in real quantum theory is efficient (Section~7.1). & Algorithmic & Yes but out
     of scope (needs a real-Hilbert-space simulator) & \textbf{No} \\
\bottomrule
\end{longtable}

\section{Method}

\subsection{Environment}
\begin{itemize}
\item Python 3.14.6, NumPy 2.4.3, SciPy 1.18.0 (available; not needed for the reproduction);
      PyMuPDF (\texttt{fitz}) 1.27.2.3 (extraction surrogate).
\item Host: CherryRd (macOS Darwin 25.3.0, x86\_64). Runtime: $<$1~s wall.
\item Statevector linear algebra only. No external quantum SDK. No LLM used for the reproduction itself.
\item Random seed: \texttt{20260705}, PRNG: \texttt{numpy.random.default\_rng}.
\end{itemize}

\subsection{Construction (concrete instantiation of Diagram 5.2)}
\begin{enumerate}
\item \textbf{Register layout.} Register~1 is the group register $\mathcal{H}_G = \mathbb{C}^{|G|}$
      with the computational basis $\{|g\rangle : g \in G\}$; register~2 is
      $\mathcal{H}_Y = \mathbb{C}^{2^M}$ where $M = \lceil \log_2 |G/H| \rceil$.
\item \textbf{Hiding function.} $f: G \to \mathbb{Z}_2^M$ constant on cosets of $H$ and injective
      on $G/H$. We assign each of the $|G/H|$ cosets a distinct random label in $\{0,\ldots,2^M-1\}$
      via \texttt{rng.choice(2**M, size=|G/H|, replace=False)}.
\item \textbf{Initial state.} $|\psi_0\rangle = |0\rangle_G \otimes |0\rangle_Y$;
      apply uniform-superposition prep on register~1 (equivalently $\mathrm{QFT}_G$ applied to $|0\rangle_G$
      for cyclic $G$):
      $|\psi_1\rangle = \tfrac{1}{\sqrt{|G|}} \sum_g |g\rangle |0\rangle_Y$.
\item \textbf{Oracle.} $U_f |g\rangle |y\rangle = |g\rangle |y \oplus f(g)\rangle$, giving
      $|\psi_2\rangle = \tfrac{1}{\sqrt{|G|}} \sum_g |g\rangle |f(g)\rangle$.
\item \textbf{Character-basis measurement of register 1.} Apply $\mathrm{QFT}_N$ to register 1
      (this is the concrete Hilbert-space form of ``measure in the character basis'' of Diagram 5.2)
      using $\mathrm{QFT}[g,y] = \omega^{gy}/\sqrt N$, $\omega = e^{2\pi i / N}$. Then measure
      in the computational basis.
\item \textbf{Joint distribution.} We compute the entire joint $|A[y,b]|^2 = |\langle y | \mathrm{QFT}\, \psi_2[:,b]\rangle|^2$
      \emph{analytically} (no shot sampling) --- this is a $|G| \times 2^M$ real matrix summing to 1.
\end{enumerate}

\subsection{Verifications}
\begin{itemize}
\item \textbf{V1: character marginal.} $P(y) = \sum_b |A[y,b]|^2$; compare to analytic
      $\mathrm{uniform}(H^\perp)$ under $L^2$ norm.
\item \textbf{V2: per-coset conditionals.} For each $b$ with $P(b)>0$, form
      $P(y|b) = |A[y,b]|^2 / P(b)$ and check (i) it is $\mathrm{uniform}(H^\perp)$,
      (ii) every $b$ gives the \emph{same} conditional (this is Diagram 5.3's
      $b$-independence claim).
\item \textbf{V3: paper's Sec~5.7 rewrite.} Pipeline~A: full protocol as above, read $P(y)$.
      Pipeline~B: $\rho_G = \mathrm{Tr}_Y |\psi_2\rangle\langle\psi_2|$, then $P(y) = \mathrm{diag}(\mathrm{QFT}\, \rho_G\, \mathrm{QFT}^\dagger)$
      (the mixed-state form obtained by applying $s^\dagger s = \mathrm{id}$ and summing over
      $b \in \mathrm{im}(s)$). Verify $\|P_A - P_B\|_\infty < 10^{-10}$ on 5 random
      hiding functions per test group.
\end{itemize}

\subsection{Commands}
\begin{lstlisting}
cd report/evidence
python3 hsp_abelian.py --seed 20260705 --outdir . > hsp_output.log
\end{lstlisting}
Full JSON result: \texttt{report/evidence/hsp\_results.json}.

\section{Results}

\subsection{Case (a): $G = \mathbb{Z}_8$, $H = \langle 2\rangle = \{0,2,4,6\}$}
Analytic $H^\perp = \{0, 4\}$; $|H|=4$, $|G/H|=2$, $M=1$.
\begin{center}
\begin{tabular}{rlrl}
\toprule
$y$ & $P(y)$ & $y$ & $P(y)$ \\
\midrule
0 & $0.5000000000$ (in $H^\perp$) & 4 & $0.5000000000$ (in $H^\perp$) \\
1 & $0$                            & 5 & $0$ \\
2 & $0$                            & 6 & $0$ \\
3 & $0$                            & 7 & $0$ \\
\bottomrule
\end{tabular}
\end{center}
\begin{itemize}
\item $P(y \in H^\perp) = 1.000000000000000$
\item $\|P - \mathrm{uniform}(H^\perp)\|_2 = 3.14 \times 10^{-16}$
\item Both cosets $\{0,2,4,6\}$, $\{1,3,5,7\}$ yield the same conditional (Diagram 5.3).
\item V3 (Sec 5.7 rewrite): max$|P_A - P_B| = 1.11 \times 10^{-16}$ across 5 random hiding functions.
\end{itemize}

\subsection{Case (b): $G = \mathbb{Z}_{15}$, $H = \langle 5\rangle = \{0,5,10\}$}
Analytic $H^\perp = \{0, 3, 6, 9, 12\}$; $|H|=3$, $|G/H|=5$, $M=3$.
\begin{center}
\begin{tabular}{rlrlrl}
\toprule
$y$ & $P(y)$ & $y$ & $P(y)$ & $y$ & $P(y)$ \\
\midrule
0  & $0.2$ (in $H^\perp$) & 5 & $0$                     & 10 & $0$ \\
1  & $0$                  & 6 & $0.2$ (in $H^\perp$)    & 11 & $0$ \\
2  & $0$                  & 7 & $0$                     & 12 & $0.2$ (in $H^\perp$) \\
3  & $0.2$ (in $H^\perp$) & 8 & $0$                     & 13 & $0$ \\
4  & $0$                  & 9 & $0.2$ (in $H^\perp$)    & 14 & $0$ \\
\bottomrule
\end{tabular}
\end{center}
\begin{itemize}
\item $P(y \in H^\perp) = 0.999999999999995$ (last digit rounding of $5 \times 0.2$)
\item $\|P - \mathrm{uniform}(H^\perp)\|_2 = 2.69 \times 10^{-15}$
\item All 5 cosets $\{0,5,10\}, \{1,6,11\}, \{2,7,12\}, \{3,8,13\}, \{4,9,14\}$
      yield the same conditional (Diagram 5.3).
\item V3 (Sec 5.7 rewrite): max$|P_A - P_B| = 2.78 \times 10^{-17}$ across 5 random hiding functions.
\end{itemize}

\section{Per-claim: what worked / what didn't}

\textbf{C1 (uniform sampling of $\mathrm{Ann}[H]$, Diagram 5.3):} \textbf{Worked.}
Both test groups produced the analytic $H^\perp$ exactly on-the-nose; $L^2$ deviation
$< 3\times 10^{-15}$; per-coset conditionals identical; probability of sampling outside
$H^\perp$ is $\le 5\times 10^{-15}$ (machine noise).

\textbf{C2 (paper's algebraic rewrites, Sec 5.3--5.8):} \textbf{Corollary worked.}
We verified V3 (the key isometry-cancellation rewrite of Eq 5.7) directly: replacing the
``measure $b$, project, then Fourier'' pipeline with the ``trace out register 2 into a mixed
state, then Fourier'' pipeline (which is what applying $s^\dagger s = \mathrm{id}$ and summing
over $b \in \mathrm{im}(s)$ delivers) gives the \emph{same} output distribution to $\approx 10^{-16}$
across 10 random hiding functions. Formal execution of the full rewrite chain requires a ZX
engine (e.g.\ PyZX \texttt{full\_reduce}); we did not run it.

\textbf{C3 (extension beyond fdHilb):} \textbf{Not tested.}
The paper's proof is category-theoretic and applies to real quantum theory, some infinite
abelian groups, and non-abelian groups with strong hypotheses. Numerical statevector
simulation is not the right tool.

\textbf{C4 (Simon in real quantum theory):} \textbf{Not tested.}
This would require a real-Hilbert-space simulator; standard numpy is complex-only.

\section{Verdict}

\textbf{REPLICATED.} The abelian HSP quantum subroutine (the concrete Hilbert-space content of
the paper's diagrammatic result, Sections 4--5, Diagram 5.3) was empirically reproduced with
real statevector simulation on two nontrivial test groups: $G=\mathbb{Z}_8, H=\langle 2\rangle$
and $G=\mathbb{Z}_{15}, H=\langle 5\rangle$. In both cases the sampled character lies in the
analytic annihilator $H^\perp$ with probability 1, uniformly, and independently of the register-2
outcome. The paper's key algebraic rewrite of Section 5.7 (the $s^\dagger s = \mathrm{id}$
isometry cancellation) was verified to preserve the register-1 marginal exactly on 10 random
hiding functions per test group. All numerical deviations from the analytic prediction are
$\lesssim 10^{-15}$, i.e.\ within floating-point noise.

\section{Open Questions}

\textbf{Q1.} \emph{Does the paper's non-abelian HSP extension (Section~6, requiring a strongly
complementary pair with a normal-classical-states property on the group observable) yield a
concrete Hilbert-space instantiation for any non-abelian group of interest --- e.g.\ the
dihedral group $D_n$ (relevant to lattice cryptography via Regev's reduction)?}
The paper only shows the abstract diagrammatic proof goes through with the extra assumption
that all classical states of the group observable are normal; it does not exhibit a Hilbert-space
example. If no non-trivial non-abelian instance exists, the extension is vacuous, which would
be an important negative result to make explicit.

\textbf{Q2.} \emph{For our two test cases, we chose the hiding function labels uniformly at
random over $\{0,\ldots,2^M-1\}$ with rejection. Does the (in)equality of the register-1
marginal from a shot-based simulation vs the analytic one degrade for certain adversarial
hiding-function label choices, or does the paper's proof imply exact-uniformity holds for
\emph{every} valid $f$?} Our V3 check across 5 random $f$s per group returned max deviation
$<3\times 10^{-16}$, but a systematic sweep over all valid label permutations of $G/H$ into
$\mathbb{Z}_2^M$ would confirm/refute label-choice independence and could reveal any
hidden numerical instability in the QFT step for large $|G|$.

\textbf{Q3.} \emph{Can the paper's Section~5.7 isometry-cancellation rewrite ($s^\dagger s = \mathrm{id}$)
be verified \emph{diagrammatically} (not just numerically) using PyZX or QuiZX \texttt{full\_reduce},
and if so, does the rewrite reach the target normal form (Diagram 5.3 RHS) within a bounded
number of steps for every abelian $G$ we tested?} Our V3 check is a numerical corollary; a
PyZX-level check would confirm that the paper's proof pipeline compiles into a mechanical
ZX-rewrite sequence for the ordinary $\dagger$-SCFA / group-algebra pair encoding of
$\mathbb{Z}_8$ and $\mathbb{Z}_{15}$.

\textbf{Q4.} \emph{The paper's Diagram 5.3 asserts $P(b, \chi) = |H|^2/|G|^2$ for
$(b, \chi) \in \mathrm{im}(s) \times \mathrm{Ann}[H]$. We verified this indirectly via the
character marginal; how does the joint $(b,y)$ distribution behave under a partial trace on
either register, and in particular does the register-2 marginal $P(b)$ exactly recover
$|H|/|G|$ per coset in our runs?} We computed the joint matrix $|A[y,b]|^2$ but only
reported the register-1 marginal and per-coset conditionals; a full statistical audit of
the joint distribution against Diagram 5.3 is a natural extension.

\textbf{Q5.} \emph{The concrete instantiation of the paper's ``group Fourier transform on
register 1'' step used the standard $\mathrm{QFT}_N$ with sign convention
$\omega = e^{+2\pi i /N}$. The character annihilator $H^\perp$ for $\mathbb{Z}_N$ is symmetric
under sign flip, so both conventions give identical support. Does the paper's diagrammatic
proof \emph{fix} a sign convention (via the antipode of the group algebra), or is it
convention-independent, and does the choice matter for non-cyclic products
$\mathbb{Z}_{n_1} \times \mathbb{Z}_{n_2}$ where $H^\perp$ may not be sign-symmetric?}

Machine-readable form in \texttt{report/open\_questions.json}.

\section*{Artifacts}
See \texttt{report/artifacts\_summary.md}.

\end{document}
